\subsection{Hiệu năng truy vấn}
Nhóm đo thời gian của bốn truy vấn phân tích trên dữ liệu đầy đủ (112.650 dòng hàng) bằng EXPLAIN ANALYZE, lấy trung bình năm lần. Trên các truy vấn quét toàn bộ dữ liệu, lược đồ hình sao chỉ nhanh hơn lược đồ chuẩn hoá của hệ thống nguồn từ 1,0 đến 1,6 lần vì dữ liệu còn nhỏ; lợi ích chính của mô hình chiều là truy vấn đơn giản và nhất quán. Chỉ mục trên khoá ngoại không thay đổi các truy vấn quét toàn bộ nhưng làm truy vấn chọn lọc theo một seller nhanh gấp khoảng chín lần. Đáng kể nhất là materialized view: đọc kết quả tổng hợp sẵn mất dưới một mili giây so với 150–750 mili giây khi view thường phải tính lại, với chi phí là phải làm mới mỗi lần nạp dữ liệu (khoảng hai giây). Fact snapshot định kỳ theo seller-tháng cũng làm scorecard seller nhanh gấp hai lần so với tính từ dòng hàng.

\begin{table}[htbp]
\centering\small
\caption{Thời gian thực thi trung bình của 4 truy vấn (mili giây, trung bình 5 lần EXPLAIN ANALYZE)}
\label{tab:benchmark}
\begin{tabular}{p{4.6cm}rrrrr}
\toprule
\textbf{Truy vấn} & \textbf{oltp (chuẩn hoá)} & \textbf{dw (star, có index FK)} & \textbf{dw không index FK} & \textbf{View thường} & \textbf{Materialized view} \\
\midrule
Q1 GMV theo bang × tháng & 657.4 & 506.1 & 494.0 & 770.4 & 0.056 \\
Q2 Lead time theo bang & 85.7 & 53.8 & 51.5 & 165.8 & 0.042 \\
Q3 Scorecard seller & 375.1 & 386.8 & 391.3 & 196.3 & 0.193 \\
Q4 GMV của một seller (tra cứu chọn lọc) & 6.3 & 0.8 & 7.3 & -- & -- \\
\bottomrule
\end{tabular}
\end{table}

\hinh{3-5-1_benchmark.png}{Thời gian thực thi trung bình: oltp so với dw, có và không có index khóa ngoại, view thường so với materialized view (thang logarit)}{fig:benchmark}{1.0}
\paragraph{Nhận xét.} Trên dữ liệu khoảng 100 nghìn đơn, lược đồ hình sao chỉ nhanh hơn lược đồ chuẩn hóa của hệ thống nguồn từ 1,0 đến 1,6 lần ở các truy vấn quét toàn bộ, và scorecard seller không nhanh hơn; lợi ích chính của mô hình chiều là \textbf{truy vấn đơn giản và số liệu nhất quán} (một định nghĩa cho mỗi KPI). Tốc độ đến chủ yếu từ hai kỹ thuật: \textbf{materialized view} (đọc dưới 1 ms thay vì 150--750 ms) và \textbf{chỉ mục} trên khóa ngoại (truy vấn chọn lọc một seller nhanh gấp khoảng 9 lần). Cần lưu ý đo trên một máy, một người dùng, dữ liệu còn nhỏ.
\todo{Nhóm xác nhận số đo trên máy của mình nếu chạy lại (python scripts/benchmark.py); thêm giải thích kế hoạch thực thi (EXPLAIN) cho một truy vấn nếu muốn.}
